% 08-bias-and-mae. The three remaining rasters on one frame, at David's
% direction. All three are MERRA-2 minus MODSCAG for water year 2023 — the
% secondary, single-year satellite validation, not the SPIReS composite on 07.
%
% Layout follows the figures' own shapes rather than a grid: the monthly panel
% grid is portrait (1453 x 1690, aspect 0.860) and takes the left column at
% full height; the two landscape figures (2.465 and 1.978) stack in the right.
%
% SIZES ARE SET AS HEIGHTS, all three, and the two on the right must sum with
% their gap to the one on the left — 2.75 + 0.26 + 3.44 = 6.45 against 6.45.
% Widths follow from the aspects (5.55, 6.78, 6.80 cm) and all clear their
% columns. Height, not width, is what binds this frame: the three fit inside
% 14.9cm for any height up to 7.5cm, and the body has about 6.5. Sizing the right pair by \linewidth instead is what overran the
% page by 62pt on the first build: two wide figures at full column width need
% about 0.91x that width in stacked height, which is more than the body has.
% Re-check the sum against the left height if any of the three changes.
\begin{frame}[t]
\frametitle{fSCA Bias and MAE, Water Year 2023}
\vspace{0.35em}
\begin{columns}[c,onlytextwidth]
\begin{column}{0.385\textwidth}
{\centering\includegraphics[height=6.45cm]{water_year_2023_monthly_spatial_grid.png}\par}
\end{column}
\begin{column}{0.600\textwidth}
{\centering
\includegraphics[height=2.75cm]{water_year_2023_bias_mae_spatial.png}\\[0.26cm]
\includegraphics[height=3.44cm]{water_year_2023_monthly_bias_mae.png}\par}
\end{column}
\end{columns}
\vspace{0.30em}
\figcaption{MERRA-2 runs low against the satellite through accumulation and
converges at melt-out.}
\end{frame}
